\section{Introduction}
Full network pre-training \citep{dai2015semi,radford2018improving, devlin2018bert,howard2018universal} has led to a series of breakthroughs in language representation learning. Many nontrivial NLP tasks, including those that have limited training data, have greatly benefited from these pre-trained models. One of the most compelling signs of these breakthroughs is the evolution of machine performance on a reading comprehension task designed for middle and high-school English exams in China, the RACE test~\citep{lai2017race}: the paper that originally describes the task and formulates the modeling challenge reports then state-of-the-art machine accuracy at $44.1\%$; the latest published result reports their model performance at $83.2\%$~\citep{liu2019roberta}; the work we present here pushes it even higher to $89.4\%$, a stunning $45.3\%$ improvement that is mainly attributable to our current ability to build high-performance pretrained language representations.

%{\em generic}\footnote{We call them generic because they improve performance not only on a specific task like RACE, but over a large battery of both language understanding and language generation tasks.}



Evidence from these improvements reveals that a large network is of crucial importance for achieving state-of-the-art performance \citep{devlin2018bert, radford2019language}. It has become common practice to pre-train large models and distill them down to smaller ones \citep{sun2019patient, turc2019well} for real applications. Given the importance of model size, we ask: \textit{Is having better NLP models as easy as having larger models}? 


%that larger pre-trained models often lead to better performance.



% A major lesson from these improvements is having large pre-trained models \citep{devlin2018bert, radford2019language} and distilling them down to smaller ones \citep{sun2019patient, turc2019well} for real applications.
% Evidence shows that larger pre-trained models often lead to better performance~\citep{devlin2018bert, radford2019language}.
% Given the importance of model size, we ask: \textit{Is having better NLP models as easy as having larger models}? 


% Past work has shown that increasingly larger models improve performance on downstream tasks~\citep{devlin2018bert, radford2018improving}.
% However, our answer to the question ``Is getting better NLU performance as simple as increasing the capacity\footnotemark[1] of the current models?'' is ``unlikely'', as supported by the following observation: when keeping the optimization procedure fixed (number of optimization steps and learning rate), doubling the hidden size of the BERT-large model leads to worse performance.
% This is illustrated by the results in Table~\ref{table:large_vs_extra_large} and the graphs in Fig.~\ref{fig:large_vs_extra_large}.
% Given that we see no evidence of overfitting in Fig.~\ref{fig:large_vs_extra_large}, we hypothesize that this is similar to the model degradation problem shown by~\cite{resnet2016}, where they conclude that not all models are similarly easy to optimize.

% This leads us to our second question: ``Is the existing capacity\footnotemark[1] of current BERT models used most efficiently?''
% The answer here is a definite ``no''.
% We show that our alternative model formulation, which we call ALBERT -- A Lite BERT --
% achieves similar or better performance while having lower model capacity, and achieves significantly better performance as its capacity increases.

% \footnotetext[1]{While it is unclear how to best compare model capacity across neural architectures, number of parameters is a useful coarse measure of capacity, especially if the architectures are similar.}

% \begin{figure}[!htbp] 
% \begin{subfigure}{.5\textwidth}
%   \centering
%   \includegraphics[width=.8\linewidth]{large_vs_extra_large_train}
%   \label{fig:large_vs_extra_large_train}
% \end{subfigure}%
% \begin{subfigure}{.5\textwidth}
%   \centering
%   \includegraphics[width=.8\linewidth]{large_vs_extra_large_dev_accuracy}
%   \label{fig:large_vs_extra_large_dev}
% \end{subfigure}
% \caption[[Caption for LOF]{Training loss (left) and dev masked LM accuracy (right) of BERT-large and BERT-xlarge (2x larger than BERT-large in terms of hidden size). The larger model has lower masked LM accuracy while showing no obvious sign of over-fitting.}
% \label{fig:large_vs_extra_large}
% \end{figure}


An obstacle to answering this question is the memory limitations of available hardware.
Given that current state-of-the-art models often have hundreds of millions or even billions of parameters, it is easy to hit these limitations as we try to scale our models.
Training speed can also be significantly hampered in distributed training, as the communication overhead is directly proportional to the number of parameters in the model.
% We also observe that simply growing the hidden size of a model such as \bertsize{large} \citep{devlin2018bert} can lead to worse performance.
% Table \ref{table:large_vs_extra_large} and Fig.~\ref{fig:large_vs_extra_large} show a typical example, where we simply increase the hidden size of BERT-large to be 2x larger and get worse results with this BERT-xlarge model. 


% \begin{table}[!htbp] 
% \small
% \centering
% \begin{tabular}{l   | c c c}
% Model  & Hidden Size & Parameters	 & RACE (Accuracy)	 \\
% \hline
%     \bertsize{large} \citep{devlin2018bert}      &1024 &334M    &72.0\% \\
%     \hline 
%     \bertsize{large} (ours)	        & 1024                          &334M	 &73.9\%\\
%     \bertsize{xlarge} (ours)      & 2048              	&1270M	 &54.3\%\\
%     %\hline
% \end{tabular}
% \caption{Increasing hidden size of BERT-large leads to worse performance on RACE.}
% \label{table:large_vs_extra_large}
% \vspace{-5pt}
% \end{table}



%\footnotetext{A sliding window smoothing is used on training loss to reduce noise from evaluating different batches.}


Existing solutions to the aforementioned problems include model parallelization \citep{shazeer2018mesh,shoeybi2019megatronlm} and clever memory management \citep{chen2016training,gomez2017reversible}.
These solutions address the memory limitation problem, but not the communication overhead.
In this paper, we address all of the aforementioned problems, by designing A Lite BERT (ALBERT) architecture that has significantly fewer parameters than a traditional BERT architecture. 

ALBERT incorporates two parameter reduction techniques that lift the major obstacles in scaling pre-trained models.
The first one is a factorized embedding parameterization.
By decomposing the large vocabulary embedding matrix into two small matrices, we separate the size of the hidden layers from the size of vocabulary embedding.
This separation makes it easier to grow the hidden size without significantly increasing the parameter size of the vocabulary embeddings.
The second technique is cross-layer parameter sharing.
This technique prevents the parameter from growing with the depth of the network.
Both techniques significantly reduce the number of parameters for BERT without seriously hurting performance, thus improving  parameter-efficiency.
An ALBERT configuration similar to BERT-large has 18x fewer parameters and can be trained about 1.7x faster.
The parameter reduction techniques also act as a form of regularization that stabilizes the training and helps with generalization. 

To further improve the performance of ALBERT, we also introduce a self-supervised loss for sentence-order prediction (SOP). SOP primary focuses on inter-sentence coherence and is designed to address the ineffectiveness \citep{yang2019xlnet, liu2019roberta} of the next sentence prediction (NSP) loss proposed in the original \bert. 


As a result of these design decisions, we are able to scale up to much larger ALBERT configurations that still have fewer parameters than BERT-large but achieve significantly better performance.
We establish new state-of-the-art results on the well-known GLUE, SQuAD, and RACE benchmarks for natural language understanding.
Specifically, we push the RACE accuracy to $89.4\%$, the GLUE benchmark to 89.4,  and the F1 score of SQuAD 2.0 to 92.2.

% capacity reallocation \footnote{While it is unclear how to best compare model capacity across neural architectures, number of parameters is a useful coarse measure of capacity, especially if the architectures are similar.}

% of the parameter reduction techniques and the new self-supervised loss